\documentclass[10pt]{article}
\usepackage[a4paper,margin=16mm]{geometry}
\usepackage[no-math]{fontspec}
\setmainfont{Latin Modern Roman}
\usepackage{amsmath,amsfonts,amssymb,amsthm,graphicx,enumitem,xurl}
\usepackage[unicode,hidelinks]{hyperref}
\newtheorem{theorem}{Theorem}[section]
\newtheorem{lemma}[theorem]{Lemma}
\newtheorem{proposition}[theorem]{Proposition}
\newtheorem{definition}[theorem]{Definition}
\newtheorem{corollary}[theorem]{Corollary}
\newtheorem{remark}[theorem]{Remark}
\allowdisplaybreaks
\setlength\emergencystretch{3em}
  \renewenvironment{proof}[1][]{
    	\begin{trivlist}
     	\item[\hspace{\labelsep}{\em\noindent Proof#1:\/}]}
     	{{\hfill$\Box$}
    	\end{trivlist}
  }
  \DeclareMathOperator{\var}{Var}
  \renewcommand{\Pr}{\mbox{\rm Pr}}
  \newcommand{\Exp}{{\mathbb{E}}}
  \newcommand{\E}{\mathbb{E}}
  \DeclareMathOperator{\gl}{GL}
  \newcommand{\R}{\mathbb{R}} 
  \newcommand{\C}{\mathbb{C}} 
  \newcommand{\D}{\mathbb{D}} 
  \newcommand{\N}{\mathbb{N}} 
  \newcommand{\Z}{\mathbb{Z}} 
  \newcommand{\F}{\mathbb{F}} 
  \newcommand{\K}{\mathbb K} 
  \newcommand{\T}{\mathbb T} 
  \newcommand{\pmset}[1]{\{-1,1\}^{#1}} 
  \newcommand{\bset}[1]{\{0,1\}^{#1}} 
  \newcommand{\sphere}[1]{S^{#1-1}} 
  \newcommand{\ball}[1]{B_{#1}} 
  \newcommand{\Orth}[1]{O(\R^{#1})} 
  \DeclareMathOperator{\im}{im} 
  \DeclareMathOperator{\vspan}{Span} 
  \newcommand{\1}{\mathbf{1}}
  \DeclareMathOperator{\gram}{Gram}
  \DeclareMathOperator{\herm}{Herm}
  \DeclareMathOperator{\conv}{Conv}
  \newcommand{\Id}{\ensuremath{\mathop{\rm I}\nolimits}}
 \newcommand{\Zp}{\mathbb Z_p} 
  \DeclareMathOperator{\charac}{char}
  \newcommand{\one}{\mathbf{1}}
  \DeclareMathOperator{\supp}{supp}
  \DeclareMathOperator{\ent}{\mathcal N}
  \DeclareMathOperator{\diam}{diam}
  \newcommand{\st}{:\,} 
  \newcommand{\ie}{{i.e.}}  
  \newcommand{\eg}{{e.g.}} 
  \newcommand{\eps}{\varepsilon}
  \newcommand{\ip}[1]{\langle #1 \rangle}
  \newcommand{\cF}{{\mathcal F}}   
  \DeclareMathOperator{\U}{\mathcal{U}}
  \DeclareMathOperator{\sign}{sign}
  \newcommand{\poly}{\mbox{\rm poly}}
  \newcommand{\ceil}[1]{\lceil{#1}\rceil}
  \newcommand{\floor}[1]{\lfloor{#1}\rfloor}
  \DeclareMathOperator{\Tr}{Tr}
  \DeclareMathOperator{\diag}{diag}
  \DeclareMathOperator{\Diag}{Diag}
  \DeclareMathOperator{\spec}{Spec}
  \DeclareMathOperator{\Fchar}{char}
  \DeclareMathOperator{\enc}{\sf Enc}
  \DeclareMathOperator{\dec}{\sf Dec}
  \DeclareMathOperator{\Norm}{N}
  \DeclareMathOperator{\rank}{rk}
  \DeclareMathOperator{\cay}{Cay}
  \DeclareMathOperator{\opt}{OPT}
  \DeclareMathOperator{\sdp}{SDP}
  \DeclareMathOperator{\cut}{cut}
  \newcommand{\id}{I}
  \newcommand{\A}{\mathcal A}
  \newcommand{\B}{\mathcal B}
  \newcommand{\CC}{\mathcal C}
  \newcommand{\HS}{\mathcal H}
  \DeclareMathOperator{\symten}{Sym}
  \DeclareMathOperator{\arank}{arank}
  \DeclareMathOperator{\prank}{prank} 
  \DeclareMathOperator{\srank}{srank} 
  \DeclareMathOperator{\GR}{GR} 
  \DeclareMathOperator{\bias}{bias}
  \DeclareMathOperator{\Pol}{Pol}
  \DeclareMathOperator{\codim}{codim}
  \DeclareMathOperator{\smrank}{schmrank}
  \newcommand{\beq}{\begin{equation}}
  \newcommand{\eeq}{\end{equation}}
  \newcommand{\beqn}{\begin{equation*}}
  \newcommand{\eeqn}{\end{equation*}}
  \newcommand{\beqr}{\begin{eqnarray}}
  \newcommand{\eeqr}{\end{eqnarray}}
  \newcommand{\beqrn}{\begin{eqnarray*}}
  \newcommand{\eeqrn}{\end{eqnarray*}}
  \newcommand{\bmline}{\begin{multline}}
  \newcommand{\emline}{\end{multline}}
  \newcommand{\bmlinen}{\begin{multline*}}
  \newcommand{\emlinen}{\end{multline*}}
  \newcommand{\nc}{{\normalfont NC}}
    \newcommand{\ncz}{{\normalfont NC}$^0$}
  \newcommand{\qnc}{{\normalfont QNC}}
  \newcommand{\ac}{{\normalfont  AC}}
    \newcommand{\acz}{{\normalfont  AC}$^0$}
  \newcommand{\acc}{{\normalfont  ACC}}
  \newcommand{\qac}{{\normalfont  QAC}}
  \newcommand{\qacc}{{\normalfont  QACC}}
\begin{document}
\begin{center}\Large Problem 2531: Random tensor restrictions\end{center}
\noindent\textbf{Authors:} Jop Bri\"et and Davi Castro-Silva\par
\noindent\textbf{Work:} Random Restrictions of High-Rank Tensors and Polynomial Maps. Discrete Analysis 2024:9, DOI 10.19086/da.124610, exact arXiv 2212.13728v2 (5 November 2024), 23 physical pages.\par
\noindent\textbf{Native source and grant:} \url{https://arxiv.org/src/2212.13728v2}, CC BY4.0, \url{https://creativecommons.org/licenses/by/4.0/}. Original PDF companion separately CC BY3.0: \url{http://creativecommons.org/licenses/by/3.0/}.\par
\noindent\textbf{Difficulty:} H1: Even with Talagrand's two-point inequality available, the task requires constructing a lower-mass event through a random ordered partition and a two-witness covering argument, then propagating the resulting exponential concentration through conditional restrictions in every tensor factor. This quantitative field-independent rank argument is beyond ordinary qualifying-exam probability or multilinear algebra; the entire authored concentration and factor iteration is retained.\par
\noindent\textbf{Editorial selection and prerequisites:} Theorem 1.4 only, with full arbitrary-field and all d/sigma quantifiers. Lemma 2.5 and all concentration support are uncounted source-local core; symmetric restrictions, matrix rank and rank examples are uncounted. Named prerequisite: Talagrand 1995 Theorem 3.1.1, exactly stated as source Theorem 2.4. Schechtman 2003 Corollary 12 is credited as inspiration, not substituted for the retained proof.\par
\noindent\textbf{Changes:} Independent standard article/theorem wrapper and original counter restoration; published class, original preamble and unprinted comments excluded. Printed mathematical tokens retained; no proof/formula replacement. Original contents list is navigation only, not a version pin. No live-corpus admission or publication claimed.\par\medskip\hrule\medskip
\setcounter{section}{1}\setcounter{theorem}{1}
Since our results are independent of the field considered (which can be finite or infinite), we will always denote it by $\F$ and suppress statements of the form ``let $\F$ be a field'' or ``for every field $\F$''.
We begin by considering the case of tensors.

\begin{definition}[Tensors]\label{def:tensors}
For finite sets $X_1,\dots,X_d \subset \N$, a $d$-tensor is a map $T:X_1\times\cdots\times X_d\to \F$.
We will associate with any $d$-tensor~$T$ a multilinear map $\F^{X_1}\times\cdots\times \F^{X_d}\to \F$ and an element of $\F^{X_1}\otimes \cdots\otimes \F^{X_d}$ in the obvious way, and also denote these objects by~$T$.
\end{definition}

\begin{definition}[Restriction of tensors]
For a tensor~$T$ as in Definition~\ref{def:tensors} and subsets $I_1\subseteq X_1,\dots,I_d\subseteq X_d$, denote $I_{[d]} = I_1\times \cdots\times I_d$ and write~$T_{|I_{[d]}}$ for the restriction of~$T$ to~$I_{[d]}$.
If~$T$ is viewed as an element of $\F^{X_1}\otimes \cdots\otimes \F^{X_d}$, then~$T_{|I_{[d]}}$ is simply a sub-tensor.
\end{definition}

We define $(\F^\infty)^{\otimes d}$ to be the set of $d$-tensors over~$\F$ with finite support.
Note that the tensors defined on finite sets naturally embed into this set, and that the rank functions for tensors discussed above are invariant under this embedding.
Our main result regarding tensors is then as follows:

\begin{theorem}\label{thm:random_tensor}
For every~$d\in \N$ and~$\sigma\in (0,1]$, there exist constants $C, \kappa > 0$ such that the following holds.
For every natural rank function $\rank: (\F^{\infty})^{\otimes d} \to \R_+$ and every $d$-tensor $T \in \bigotimes_{i=1}^d \F^{[n_i]}$ we have
\beqn
\Pr_{I_1\sim [n_1]_{\sigma}, \dots, I_d\sim [n_d]_{\sigma}}\big[ \rank \big(T_{|I_{[d]}}\big) \geq \kappa\cdot \rank(T)\big] \geq 1 - C e^{-\kappa \rank(T)}.
\eeqn
\end{theorem}
\setcounter{subsection}{2}
\subsection{Notation}
\label{sec:prelims}

Here we collect some notation that will be used throughout the paper.
Given sets  $I_1, \dots, I_d$, we write $I_{[d]} = I_1\times \cdots\times I_d$.
For a set~$J$ and some~$\sigma \in [0,1]$, we denote by~$\pi_{\sigma}^J$ the product distribution on~$\{0,1\}^J$ where each coordinate is independently set to $1$ with probability~$\sigma$, or to $0$ with probability~$1-\sigma$;
when~$J=[n]$, we write simply~$\pi_{\sigma}^n$.
We write $J_\sigma$ for the probability distribution over subsets of~$J$ where each element is present independently with probability~$\sigma$.
To denote that a random variable~$X$ is distributed according to a distribution~$\mu$, we write $X\sim \mu$.
\section{Tensors}
\label{sec:tensors}

Recall that, for a field $\F$ and integer $d\geq 2$, we denote by $(\F^\infty)^{\otimes d}$ the set of all $d$-tensors of finite support over $\F$.
All our results (including the values of the implied constants) will hold independently of the field considered, so we will always denote it by $\F$ without further comment.

The notions of tensor rank we will consider here are those called \emph{natural rank functions} as defined below:

\begin{definition}[Natural rank]\label{def:nat_rank}
We say that $\rank: (\F^{\infty})^{\otimes d} \to \R_+$ is a natural rank function if $\rank({\bf 0}) = 0$ and it satisfies the following properties:
\begin{enumerate}
    \item Sub-additivity: \\
    $\rank(T+S) \leq \rank(T) + \rank(S)$ for all~$T, S\in (\F^{\infty})^{\otimes d}$.
    \item Monotonicity under restrictions: \\
    $\rank\big(T_{|I_{[d]}}\big) \leq \rank(T)$ for all~$T\in (\F^{\infty})^{\otimes d}$ and all sets~$I_1, \dots, I_d \subset \N$.
    \item Restriction Lipschitz property: \\
    $\rank\big(T_{|J_{[d]}}\big) \leq \rank\big(T_{|I_{[d]}}\big) + \sum_{i=1}^d |J_i \setminus I_i|$
    for all~$T\in (\F^{\infty})^{\otimes d}$ and all sets~$I_1\subseteq J_1, \dots, I_d\subseteq J_d$.
    \end{enumerate}
\end{definition}
\setcounter{theorem}{2}
\subsection{Concentration for monotone sub-additive functions}

The main step in our proof of Theorem~\ref{thm:random_tensor} is a type of concentration inequality for monotone sub-additive functions on the hypercube.

In order to obtain such a result we will consider a notion of \emph{two-point distance} from sets $A\subseteq \bset{n}$, which intuitively measures how many coordinates of some given point $x$ cannot be captured by any two elements of~$A$.

\begin{definition}
Given a point $x\in \bset{n}$ and a set $A \subseteq \bset{n}$, the two-point distance between $x$ and $A$ is
\beqn
h_2(x; A) = \min_{y, z \in A} \big|\big\{i \in [n] \st x_i \neq y_i \And x_i \neq z_i\big\}\big|.
\eeqn
\end{definition}

Note that this ``distance'' $h_2(x; A)$ can be zero even if $x \notin A$;
for instance, $h_2\big(00;\, \{01, 10\}\big) = 0$.
The reason for considering this notion is that one gets much better concentration for $h_2$ than one does for the usual Hamming distance.
This is shown by the following result, which is a special case of an inequality of Talagrand~\cite[Theorem 3.1.1]{Talagrand:1995}:

\begin{theorem}\label{thm:Talagrand}
For any parameter $\sigma\in [0,1]$ and set $A \subseteq \bset{n}$, we have that
\beqn
\Pr_{x\sim \pi_\sigma^n}\big[h_2(x; A) \geq k\big] \leq 2^{-k} \pi_\sigma^n(A)^{-2} \quad \text{for all $k\geq 0$.}
\eeqn
\end{theorem}

The next lemma is the key result needed for proving our main theorem for tensors, and it deals more abstractly with monotone, sub-additive Lipschitz functions on the hypercube.
We endow~$\bset{n}$ with the usual partial order, where $x\leq y$ if the support of~$x$ is contained in the support of~$y$.
A function $f:\bset{n} \to \R$ is monotone if $f(x) \leq f(y)$ whenever $x\leq y$, and it is sub-additive if $f(x+y) \leq f(x) + f(y)$ for all $x,y\in\bset{n}$ with  disjoint supports.
Finally,~$f$ is $1$-Lipschitz if
\beqn
\big|f(x) - f(y)\big| \leq \big|\big\{i \in [n] \st x_i \neq y_i\big\}\big| \quad \text{for all $x, y \in \bset{n}$}.
\eeqn
For a string $x\in \bset{n}$ and set $I\subseteq[n]$, we let~$x_I\in \bset{n}$ be the string that equals~$x$ on the indices in~$I$ and is zero elsewhere.
The lemma that follows and its proof are inspired by an argument of Schechtman~\cite[Corollary~12]{Schechtman2003}.

\begin{lemma}[Concentration inequality] \label{lem:subadditive}
Let $f:\bset{n} \to \R_+$ be a monotone, sub-additive $1$-Lipschitz function such that $f({\bf 0}) =0$ and $f({\bf 1}) = r$.
Then,
\begin{align*}
    \Pr_{x\sim \pi_{\sigma}^n} \big[f(x) \leq \sigma r/4\big] &\leq \sqrt{\frac{3}{2\sigma}}\, 2^{-\sigma r/12} \quad &\text{if $\,0 < \sigma < 1/2$,} \\
    \Pr_{x\sim \pi_{\sigma}^n} \big[f(x) \leq r/6\big] &\leq \sqrt{2}\, 2^{-r/12} &\text{if $\,1/2 \leq \sigma \leq 1$.}
\end{align*}
\end{lemma}

\begin{proof}
We first prove the first inequality above.
Let $\sigma\in (0, 1/2)$, and consider the set
\beqn
A = \big\{ x\in \bset{n} \st f(x)\leq \sigma r/4 \big\}.
\eeqn
By Talagrand's Inequality (Theorem~\ref{thm:Talagrand}) we have
\beqn
\Pr_{x\sim \pi_{\sigma}^n}\big[ h_2(x; A)\geq \sigma r/6\big] \leq 2^{-\sigma r/6} \,\Pr_{x\sim \pi_{\sigma}^n}\big[ f(x)\leq \sigma r/4\big]^{-2},
\eeqn
so to finish the proof it suffices to show that
\beqn
\Pr_{x\sim \pi_{\sigma}^n}\big[ h_2(x; A)\geq \sigma r/6\big] \geq 2\sigma/3.
\eeqn

We claim that
\beq \label{eq:h2inclusion}
\big\{ x\in \bset{n} \st f(x)\geq 2\sigma r/3 \big\} \subseteq \big\{ x\in \bset{n} \st h_2(x; A)\geq \sigma r/6 \big\}.
\eeq
Indeed, if $h_2(x; A) < \sigma r/6$ then there are $y, z\in A$ such that
\beqn
\big|\big\{i \in [n] \st x_i \neq y_i \And x_i \neq z_i\big\}\big| < \sigma r/6.
\eeqn
Denote $I = \big\{i \in [n] \st x_i \neq y_i \And x_i \neq z_i\big\}$, $J = \big\{i \in [n] \st x_i = y_i\big\}$ and $J' = \big\{i \in [n] \st x_i \neq y_i \And x_i = z_i\big\}$;
note that these sets partition~$[n]$, and by assumption $|I| < \sigma r/6$.
We then have
\begin{align*}
    f(x) &\leq f(x_I + x_J) + f(x_{J'}) \\
    &= f(x_I + y_J) + f(z_{J'}) \\
    &\leq |I| + f(y_J) + f(z_{J'}) \\
    &\leq |I| + f(y) + f(z) \\
    &< 2\sigma r/3,
\end{align*}
so $\big\{h_2(x; A) < \sigma r/6\big\} \subseteq \big\{f(x) < 2\sigma r/3\big\}$ and inclusion \eqref{eq:h2inclusion} follows.
It thus suffices to show that $\Pr_{x\sim \pi_{\sigma}^n}\big[ f(x)\geq 2\sigma r/3\big] \geq 2\sigma/3$.

We will next prove that, for any integer $k\geq 1$, we have
\beq \label{eq:coupling}
\Pr_{x\sim \pi_{1/k}^n}\big[ f(x)\geq r/k\big] \geq 1/k.
\eeq
This is done by a simple \emph{coupling argument}, which will be important for us again later on.
Consider a uniformly random ordered $k$-partition $(I_1, \dots, I_k)$ of $[n]$;
thus the $k$ sets $I_i$ are pairwise disjoint and have union $[n]$, with each of the $k^n$ possible such $k$-tuples having the same probability.
Denoting by $\1_I$ the indicator function of set $I$, we have
\beqn
r = f(\1_{[n]}) = f(\1_{I_1} + \dots + \1_{I_k}) \leq \sum_{i=1}^k f(\1_{I_i}),
\eeqn
so $f(\1_{I_i}) \geq r/k$ holds for at least one $i\in [k]$ in \emph{every} ordered partition $(I_1, \dots, I_k)$.
By symmetry, it follows that $\Pr \big[f(\1_{I_1}) \geq r/k \big] \geq 1/k$.
Since the marginal distribution of $\1_{I_1}$ (and every other $\1_{I_i}$) is precisely $\pi_{1/k}^n$, we conclude that
\begin{equation*}
\Pr_{x\sim \pi_{1/k}^n} \big[f(x) \geq r/k\big] = \Pr \big[f(\1_{I_1}) \geq r/k \big] \geq 1/k,
\end{equation*}
as wished.

Now let $k = \lceil 1/\sigma \rceil$.
Since $0 < \sigma < 1/2$, we have $2\sigma/3 \leq 1/k \leq \sigma$, and so by monotonicity of $f$
\begin{align*}
    \Pr_{x\sim \pi_{\sigma}^n}\big[ f(x)\geq 2\sigma r/3\big] &\geq \Pr_{x\sim \pi_{1/k}^n}\big[ f(x)\geq 2\sigma r/3\big] \\
    &\geq \Pr_{x\sim \pi_{1/k}^n}\big[ f(x)\geq r/k\big].
\end{align*}
From inequality~\eqref{eq:coupling} we then conclude that
\beqn
\Pr_{x\sim \pi_{\sigma}^n}\big[ f(x)\geq 2\sigma r/3\big] \geq 1/k \geq 2\sigma/3,
\eeqn
finishing the proof of the first inequality in the statement of the lemma.

The second inequality is proven using the same arguments, now applied to the parameter $\sigma = 1/2$ and the set $A = \big\{ x\in \bset{n} \st f(x)\leq r/6 \big\}$.
By the same argument as above,
\begin{equation*}
    \big\{x\in \bset{n}\st f(x) \geq r/2\big\} \subseteq \big\{x\in \bset{n}\st h_2(x;A) \geq r/6\big\}.
\end{equation*}
Talagrand's inequality and coupling imply that
\begin{align*}
    2^{-r/6}\Pr_{x\sim\pi_{1/2}^n}\big[f(x) \leq r/6\big]^{-2} &\geq \Pr_{x\sim \pi_{1/2}^n}\big[h_2(x;A) \geq r/6\big]\\
    &\geq \Pr_{x\sim\pi_{1/2}^n}\big[f(x) \geq r/2\big]\\
    &\geq \frac{1}{2}.
\end{align*}
Rearranging then gives the result for~$\sigma =  1/2$;
we obtain slightly better bounds since in this case $1/\sigma = 2$ is an integer, and so no rounding errors occur.
By monotonicity of $f$, the same bound continues to hold for all $\sigma > 1/2$.
\end{proof}


\subsection{The Random Restriction Theorem}

It is now a simple matter to use Lemma~\ref{lem:subadditive} to prove the Random Restriction Theorem for tensors.

\begin{proof}[ of Theorem~\ref{thm:random_tensor}]
Let $\rank: (\F^{\infty})^{\otimes d} \to \R$ be a natural rank function, and $T \in \bigotimes_{i=1}^d \F^{[n_i]}$ be a tensor.
We wish to show that
\beqn
\Pr_{I_1\sim [n_1]_{\sigma}, \dots, I_d\sim [n_d]_{\sigma}}\big[ \rank \big(T_{|I_{[d]}}\big) \geq \kappa(\sigma) \cdot \rank(T)\big] \geq 1 - C(\sigma) e^{-\kappa(\sigma) \rank(T)}
\eeqn
for some well-chosen constants $C(\sigma), \kappa(\sigma) > 0$ depending only on the value of $\sigma > 0$ and the order $d$ of the tensor.
We consider here the case where $\sigma < 1/2$, as the case where $\sigma \geq 1/2$ is analogous.\footnote{This second case is also an immediate consequence of first case together with monotonicity of $\rank$ under restrictions.}

Define the function $f_1: \bset{n_1} \to \R_+$ as follows:
for $x \in \bset{n_1}$ with support $J$, $f_1(x)$ is the rank of the subtensor of $T$ whose indices of the first row belong to $J$.
In formula:
\beqn
f_1(\1_{J}) = \rank\big( T_{|J \times \prod_{j=2}^d [n_j]} \big) \quad \text{for all } J\subseteq [n_1].
\eeqn
By the definition of natural rank, $f_1$ is a monotone, sub-additive $1$-Lipschitz function with maximum value $\rank(T)$.
Since $\1_J \sim \pi_{\sigma}^{n_1}$ when $J\sim [n_1]_{\sigma}$, it follows from Lemma~\ref{lem:subadditive} that
\beq \label{eq:f1}
\Pr_{I_1\sim [n_1]_{\sigma}}\bigg[ \rank\big( T_{|I_1 \times \prod_{j=2}^d [n_j]} \big) \geq \frac{\sigma \rank(T)}{4}\bigg] \geq 1 - \sqrt{\frac{3}{2\sigma}}\, 2^{-\frac{\sigma}{12} \rank(T)}.
\eeq

For any fixed $I_1 \subseteq [n_1]$ satisfying $f_1(\1_{I_1}) \geq \sigma \rank(T)/4$, we define the function $f_2: \bset{n_2} \to \R_+$ by
\beqn
f_2(\1_{J}) = \rank\big( T_{|I_1 \times J \times \prod_{j=3}^d [n_j]} \big) \quad \text{for all } J\subseteq [n_2].
\eeqn
This function is again monotone, sub-additive and $1$-Lipschitz, and it has maximum value at least~$\sigma \rank(T)/4$.
By Lemma~\ref{lem:subadditive} we have
\beqn
\Pr_{I_2\sim [n_2]_{\sigma}}\bigg[ \rank\big( T_{|I_1 \times I_2 \times \prod_{j=3}^d [n_j]} \big) \geq \Big(\frac{\sigma}{4}\Big)^2 \rank(T)\bigg] \geq 1 - \sqrt{\frac{3}{2\sigma}}\, 2^{-\frac{\sigma}{12} \frac{\sigma}{4} \rank(T)}.
\eeqn
Since this holds whenever $f_1(\1_{I_1}) \geq \sigma \rank(T)/4$, taking the union bound together with inequality~\eqref{eq:f1} we obtain
\beqn
\Pr_{I_1\sim [n_1]_{\sigma}, I_2\sim [n_2]_{\sigma}}\bigg[ \rank\big( T_{|I_1 \times I_2 \times \prod_{j=3}^d [n_j]} \big) \geq \Big(\frac{\sigma}{4}\Big)^2 \rank(T)\bigg] \geq 1 - 2\sqrt{\frac{3}{2\sigma}}\, 2^{-\frac{\sigma}{12} \frac{\sigma}{4} \rank(T)}.
\eeqn
Proceeding in this same way for each row of the tensor, and always taking the union bound, we eventually conclude that
\beqn
\Pr_{I_1\sim [n_1]_{\sigma}, \dots, I_d\sim [n_d]_{\sigma}}\bigg[ \rank \big(T_{|I_{[d]}}\big) \geq \Big(\frac{\sigma}{4}\Big)^d \rank(T)\bigg] \geq 1 - d\sqrt{\frac{3}{2\sigma}}\, 2^{-\frac{\sigma}{12} (\frac{\sigma}{4})^{d-1} \rank(T)},
\eeqn
which finishes the proof.
\end{proof}

\begin{thebibliography}{10}

\bibitem{AnanyanH:2020}
Tigran Ananyan and Melvin Hochster.
\newblock Small subalgebras of polynomial rings and {S}tillman's conjecture.
\newblock {\em J. Amer. Math. Soc.}, 33:291--309, 2020.
\newblock \href {https://doi.org/10.1090/jams/932} {\path{doi:10.1090/jams/932}}.

\bibitem{BCCGNSU}
Jonah Blasiak, Thomas Church, Henry Cohn, Joshua~A. Grochow, Eric Naslund, William~F. Sawin, and Chris Umans.
\newblock On cap sets and the group-theoretic approach to matrix multiplication.
\newblock {\em Discrete Anal.}, pages Paper No. 3, 27, 2017.
\newblock \href {https://doi.org/10.19086/da.1245} {\path{doi:10.19086/da.1245}}.

\bibitem{BlatterDR:2022}
Andreas Blatter, Jan Draisma, and Filip Rupniewski.
\newblock A tensor restriction theorem over finite fields, 2022.
\newblock arXiv:2211.12319.
\newblock \href {https://doi.org/10.48550/ARXIV.2211.12319} {\path{doi:10.48550/ARXIV.2211.12319}}.

\bibitem{Briet:2021}
Jop Bri\"{e}t.
\newblock Subspaces of tensors with high analytic rank.
\newblock {\em Online J. Anal. Comb.}, (16):Paper No. 6, 9, 2021.

\bibitem{BrietBCSN}
Jop Briët, Harry Buhrman, Davi Castro-Silva, and Niels M.~P. Neumann.
\newblock Noisy decoding by shallow circuits with parities: classical and quantum, 2023.
\newblock \href {https://doi.org/10.48550/ARXIV.2302.02870} {\path{doi:10.48550/ARXIV.2302.02870}}.

\bibitem{CohenM21}
Alex Cohen and Guy Moshkovitz.
\newblock Structure vs.\ randomness for bilinear maps.
\newblock In {\em S{TOC} '21---{P}roceedings of the 53rd {A}nnual {ACM} {SIGACT} {S}ymposium on {T}heory of {C}omputing}, pages 800--808. ACM, New York, 2021.
\newblock \href {https://doi.org/10.1145/3406325.3451007} {\path{doi:10.1145/3406325.3451007}}.

\bibitem{CohenM21b}
Alex Cohen and Guy Moshkovitz.
\newblock Partition and analytic rank are equivalent over large fields.
\newblock {\em Duke Math. J.}, 172(12):2433--2470, 2023.
\newblock \href {https://doi.org/10.1215/00127094-2022-0086} {\path{doi:10.1215/00127094-2022-0086}}.

\bibitem{CrootLP2017}
Ernie Croot, Vsevolod~F. Lev, and P\'{e}ter~P\'{a}l Pach.
\newblock Progression-free sets in~{$\Bbb Z^n_4$} are exponentially small.
\newblock {\em Ann. of Math. (2)}, 185(1):331--337, 2017.
\newblock \href {https://doi.org/10.4007/annals.2017.185.1.7} {\path{doi:10.4007/annals.2017.185.1.7}}.

\bibitem{DerksenES:2017}
Harm Derksen, Rob~H. Eggermont, and Andrew Snowden.
\newblock Topological noetherianity for cubic polynomials.
\newblock {\em Algebra Number Theory}, 11(9):2197--2212, 2017.
\newblock \href {https://doi.org/10.2140/ant.2017.11.2197} {\path{doi:10.2140/ant.2017.11.2197}}.

\bibitem{Draisma:2019}
Jan Draisma.
\newblock Topological {N}oetherianity of polynomial functors.
\newblock {\em J. Amer. Math. Soc.}, 32(3):691--707, 2019.
\newblock \href {https://doi.org/10.1090/jams/923} {\path{doi:10.1090/jams/923}}.

\bibitem{EllenbergG2017}
Jordan~S. Ellenberg and Dion Gijswijt.
\newblock On large subsets of {$\Bbb F^n_q$} with no three-term arithmetic progression.
\newblock {\em Ann. of Math. (2)}, 185(1):339--343, 2017.
\newblock \href {https://doi.org/10.4007/annals.2017.185.1.8} {\path{doi:10.4007/annals.2017.185.1.8}}.

\bibitem{FoxL17}
Jacob Fox and L\'{a}szl\'{o}~Mikl\'{o}s Lov\'{a}sz.
\newblock A tight bound for {G}reen's arithmetic triangle removal lemma in vector spaces.
\newblock {\em Adv. Math.}, 321:287--297, 2017.
\newblock \href {https://doi.org/10.1016/j.aim.2017.09.037} {\path{doi:10.1016/j.aim.2017.09.037}}.

\bibitem{Friedgut1998}
Ehud Friedgut.
\newblock Boolean functions with low average sensitivity depend on few coordinates.
\newblock {\em Combinatorica}, 18(1):27--35, 1998.
\newblock \href {https://doi.org/10.1007/PL00009809} {\path{doi:10.1007/PL00009809}}.

\bibitem{GowersK2022}
W.~T. Gowers and Thomas Karam.
\newblock Equidistribution of high-rank polynomials with variables restricted to subsets of $\mathbb{F}_p$, 2022.
\newblock arXiv:2209.04932.
\newblock \href {https://doi.org/10.48550/ARXIV.2209.04932} {\path{doi:10.48550/ARXIV.2209.04932}}.

\bibitem{GowersW2011}
W.~T. Gowers and J.~Wolf.
\newblock Linear forms and higher-degree uniformity for functions on {$\Bbb F^n_p$}.
\newblock {\em Geom. Funct. Anal.}, 21(1):36--69, 2011.
\newblock \href {https://doi.org/10.1007/s00039-010-0106-3} {\path{doi:10.1007/s00039-010-0106-3}}.

\bibitem{GreenT2009}
Ben Green and Terence Tao.
\newblock The distribution of polynomials over finite fields, with applications to the {G}owers norms.
\newblock {\em Contrib. Discrete Math.}, 4(2):1--36, 2009.

\bibitem{HagerupRub:1990}
Torben Hagerup and Christine R{\"u}b.
\newblock A guided tour of {Chernoff} bounds.
\newblock {\em Information Processing Letters}, 33(6):305--308, 1990.
\newblock \href {https://doi.org/10.1016/0020-0190(90)90214-I} {\path{doi:10.1016/0020-0190(90)90214-I}}.

\bibitem{HaramatyS2010}
Elad Haramaty and Amir Shpilka.
\newblock On the structure of cubic and quartic polynomials.
\newblock In {\em S{TOC}'10---{P}roceedings of the 2010 {ACM} {I}nternational {S}ymposium on {T}heory of {C}omputing}, pages 331--340. ACM, New York, 2010.
\newblock \href {https://doi.org/10.1145/1806689.1806736} {\path{doi:10.1145/1806689.1806736}}.

\bibitem{Hatami2016}
Pooya Hatami.
\newblock On the structure of quintic polynomials.
\newblock In {\em Approximation, randomization, and combinatorial optimization. {A}lgorithms and techniques}, volume~60 of {\em LIPIcs. Leibniz Int. Proc. Inform.}, pages Art. No. 33, 18. Schloss Dagstuhl. Leibniz-Zent. Inform., Wadern, 2016.
\newblock \href {https://doi.org/10.4230/LIPIcs.APPROX-RANDOM.2016.33} {\path{doi:10.4230/LIPIcs.APPROX-RANDOM.2016.33}}.

\bibitem{Karam2022}
Thomas Karam.
\newblock High-rank subtensors of high-rank tensors, 2022.
\newblock arXiv:2207.08030.
\newblock \href {https://doi.org/10.48550/ARXIV.2207.08030} {\path{doi:10.48550/ARXIV.2207.08030}}.

\bibitem{KazhdanZ:cohom}
David Kazhdan and Tamar Ziegler.
\newblock Approximate cohomology.
\newblock {\em Selecta Math. (N.S.)}, 24(1):499--509, 2018.
\newblock \href {https://doi.org/10.1007/s00029-017-0335-5} {\path{doi:10.1007/s00029-017-0335-5}}.

\bibitem{KazhdanZ:2020}
David Kazhdan and Tamar Ziegler.
\newblock Properties of high rank subvarieties of affine spaces.
\newblock {\em Geom. Funct. Anal.}, 30(4):1063--1096, 2020.
\newblock \href {https://doi.org/10.1007/s00039-020-00542-4} {\path{doi:10.1007/s00039-020-00542-4}}.

\bibitem{KoppartyMZ}
Swastik Kopparty, Guy Moshkovitz, and Jeroen Zuiddam.
\newblock Geometric rank of tensors and subrank of matrix multiplication.
\newblock In {\em 35th {C}omputational {C}omplexity {C}onference}, volume 169 of {\em LIPIcs. Leibniz Int. Proc. Inform.}, pages Art. No. 35, 21. Schloss Dagstuhl. Leibniz-Zent. Inform., Wadern, 2020.
\newblock \href {https://doi.org/10.4230/LIPIcs.CCC.2020.35} {\path{doi:10.4230/LIPIcs.CCC.2020.35}}.

\bibitem{Landsberg08}
J.~M. Landsberg.
\newblock Geometry and the complexity of matrix multiplication.
\newblock {\em Bull. Amer. Math. Soc. (N.S.)}, 45(2):247--284, 2008.
\newblock \href {https://doi.org/10.1090/S0273-0979-08-01176-2} {\path{doi:10.1090/S0273-0979-08-01176-2}}.

\bibitem{Lovett2019}
Shachar Lovett.
\newblock The analytic rank of tensors and its applications.
\newblock {\em Discrete Anal.}, pages Paper No. 7, 10, 2019.
\newblock \href {https://doi.org/10.19086/da} {\path{doi:10.19086/da}}.

\bibitem{NaslundEGZ}
Eric Naslund.
\newblock Exponential bounds for the {Erd\H{o}s}-{G}inzburg-{Z}iv constant.
\newblock {\em J. Combin. Theory Ser. A}, 174:105185, 19, 2020.
\newblock \href {https://doi.org/10.1016/j.jcta.2019.105185} {\path{doi:10.1016/j.jcta.2019.105185}}.

\bibitem{Naslund2020}
Eric Naslund.
\newblock The partition rank of a tensor and {$k$}-right corners in {$\Bbb F_q^n$}.
\newblock {\em J. Combin. Theory Ser. A}, 174:105190, 25, 2020.
\newblock \href {https://doi.org/10.1016/j.jcta.2019.105190} {\path{doi:10.1016/j.jcta.2019.105190}}.

\bibitem{NaslundS17}
Eric Naslund and Will Sawin.
\newblock Upper bounds for sunflower-free sets.
\newblock {\em Forum Math. Sigma}, 5:Paper No. e15, 10, 2017.
\newblock \href {https://doi.org/10.1017/fms.2017.12} {\path{doi:10.1017/fms.2017.12}}.

\bibitem{ODonnell:2014}
Ryan O'Donnell.
\newblock {\em Analysis of {B}oolean functions}.
\newblock Cambridge University Press, New York, 2014.
\newblock \href {https://doi.org/10.1017/CBO9781139814782} {\path{doi:10.1017/CBO9781139814782}}.

\bibitem{Raz13}
Ran Raz.
\newblock Tensor-rank and lower bounds for arithmetic formulas.
\newblock {\em J. ACM}, 60(6):Art. 40, 15, 2013.
\newblock \href {https://doi.org/10.1145/2535928} {\path{doi:10.1145/2535928}}.

\bibitem{Schechtman2003}
Gideon Schechtman.
\newblock Concentration results and applications.
\newblock In {\em Handbook of the geometry of {B}anach spaces, {V}ol. 2}, pages 1603--1634. North-Holland, Amsterdam, 2003.
\newblock \href {https://doi.org/10.1016/S1874-5849(03)80044-X} {\path{doi:10.1016/S1874-5849(03)80044-X}}.

\bibitem{Schmidt:1985}
Wolfgang~M. Schmidt.
\newblock The density of integer points on homogeneous varieties.
\newblock {\em Acta Math.}, 154(3-4):243--296, 1985.
\newblock \href {https://doi.org/10.1007/BF02392473} {\path{doi:10.1007/BF02392473}}.

\bibitem{Talagrand:1995}
Michel Talagrand.
\newblock Concentration of measure and isoperimetric inequalities in product spaces.
\newblock {\em Inst. Hautes \'{E}tudes Sci. Publ. Math.}, (81):73--205, 1995.

\bibitem{Tao:srank}
Terence Tao.
\newblock \url{https://terrytao.wordpress.com/2016/05/18/a-symmetric-formulation-of-the-croot-lev-pach-ellenberg-gijswijt-capset-bound/}, 2016.

\bibitem{TaoSawin}
Terence Tao and Will Sawin.
\newblock \url{https://terrytao.wordpress.com/2016/08/24/notes-on-the-slice-rank-of-tensors/}, 2016.

\bibitem{TaoZ12}
Terence Tao and Tamar Ziegler.
\newblock The inverse conjecture for the {G}owers norm over finite fields in low characteristic.
\newblock {\em Ann. Comb.}, 16(1):121--188, 2012.
\newblock \href {https://doi.org/10.1007/s00026-011-0124-3} {\path{doi:10.1007/s00026-011-0124-3}}.

\end{thebibliography}
\end{document}
